\documentclass[11pt]{article}
\usepackage[margin=0.9in]{geometry}
\usepackage{amsmath,amssymb,microtype}
\usepackage[hidelinks]{hyperref}

\title{The Grothendieck-type exponent conjecture is already answered}
\author{Literature status for arXiv:1102.4542}
\date{13 August 2026}

\begin{document}
\maketitle

\begin{center}
\fbox{\parbox{0.91\textwidth}{\textbf{Status: literature already answered.}
Pellegrino--Seoane-Sep\'ulveda, arXiv:1307.4809, explicitly identifies and
solves Bernardino's 2011 conjecture.}}
\end{center}

\section*{Source question}

A. T. Bernardino, \emph{On cotype and a Grothendieck-type theorem for
absolutely summing multilinear operators}, arXiv:1102.4542; Quaest. Math.
\textbf{34} (2011), 513--519.

Problem 1.3 on PDF page 2 asks for the best $g_n>0$ such that
\[
 \Pi^n_{(g_n;1,\ldots,1)}(\ell_1,\ldots,\ell_1;\ell_2)
 =\mathcal L(\ell_1,\ldots,\ell_1;\ell_2).
\]
The final theorem proves $g_n\le 2/(n+1)$, and the last paragraph conjectures
that $2/(n+1)$ is optimal.

\section*{Exact later answer}

D. Pellegrino and J. B. Seoane-Sep\'ulveda,
\emph{Grothendieck's theorem for absolutely summing multilinear operators is
optimal}, arXiv:1307.4809; Linear Multilinear Algebra \textbf{63} (2015),
554--558, DOI \href{https://doi.org/10.1080/03081087.2013.877013}{10.1080/03081087.2013.877013}.

Its abstract states that the result ``solves (in the positive) a conjecture
posed by A. T. Bernardino in 2011.''  Theorem 1.1 on PDF page 1 proves
\[
 \boxed{g_n=\frac{2}{n+1}.}
\]
More strongly, for every $r<2/(n+1)$ there is a continuum-dimensional linear
space whose nonzero elements are continuous $n$-linear maps
$\ell_1^n\to\ell_2$ that are not absolutely $(r;1,\ldots,1)$-summing.

\section*{Identification}

The notation and domain/range are identical after renaming the arity $m$ in
the supporting paper to $n$ in the source.  The source supplies the upper
bound, while Theorem 1.1 supplies failure for every smaller exponent.  Thus
the source's sharp-exponent problem and final conjecture are fully resolved.
The source's broader Problem 1.5 about optimal cotype-dependent inclusion
exponents is a separate question and is not claimed here.

\end{document}
